\documentclass[10pt]{article}
\usepackage[margin=.85in]{geometry}
\usepackage{amsmath,amssymb,amsthm,booktabs}
\newcommand{\E}{\mathbb E}\newcommand{\Pp}{\mathbb P}
\newtheorem{theorem}{Theorem}
\title{A low-noise support bound on thirteen probability panels}
\author{}\date{28 October 2026}
\begin{document}\maketitle
\noindent This independent check extends the direct low-noise argument to
all thirteen closed panels in \(2/5\le p\le53/100\). It uses no Fourier
kernel, phase export or imported numerical program. The contact and
Gaussian estimates below remain explicit analytic premises.

\paragraph{Precise conditional statement.}
Let \(X\) have mean zero and variance one, with an actual central gap
\((-A,B)\), \(A,B>0\), and let \(p=\Pp(X>0)\), \(q=1-p\).
Write
\[
 m=\E X_+=\E X_-,\quad \ell=m/(pq),\quad
 X=\ell(Z-p+U),\quad Z=\mathbf1_{\{X>0\}},\quad
 \E(U\mid Z)=0 .
\]
The sign of \(Z-p+U\) agrees with that of \(Z-p\). Put
\[
 \theta=\ell^2\E U^2,\qquad \lambda=n\E U^2,\qquad
 \delta=1-(A+B)/\ell.
\]
Assume \(n\ge42\), \(0\le\lambda\le3/40\),
\(2/5\le p\le53/100\), the already acquired
\(|U|\le3/5\), and \(\delta<.072\).
Assume the fixed-target closed-support contact inequality with a sharp
predecessor and nonnegative excess, and its uniform Gaussian bounds
specified below. These premises apply to the same law.
\begin{theorem}
Under these hypotheses \(|U|<1/2\), including at every support atom.
The conclusion holds for every raw threshold \(H\) and every actual
\(n\ge42\). It makes no claim for larger noise or for other probability
panels.
\end{theorem}

\paragraph{Probability and moment bounds.}
Fix one closed interval \(I=[l,u]=[i/100,(i+1)/100]\),
\(i=40,\ldots,52\), and write
\[
 v_-=\min_{p\in I}p(1-p),\quad v_+=\max_{p\in I}p(1-p),\quad
 \Theta=\frac{3/40}{42v_-+3/40},\qquad
 E_I\le\sqrt{\frac{1-\Theta}{v_+}}.
\]
The maximum \(v_+\) includes \(1/4\) when the interval meets \(1/2\).
Variance gives \(\ell^2(pq+\lambda/n)=1\), hence
\(\theta=\lambda/(npq+\lambda)\le\Theta\) and \(\ell\ge E_I>0\).

Define \(s_+=\ell^2\E[U^2\mathbf1_{\{Z=1\}}]\) and
\(s_-=\ell^2\E[U^2\mathbf1_{\{Z=0\}}]\), so
\(s_++s_-=\theta\).
Conditional centering and sign preservation give
\[
 \beta:=\E[X|X|]=(1-\theta)(1-2p)+s_+-s_-.
\]
Both endpoint functions decrease with \(p\), including when \(p>1/2\).
Therefore
\[
 b_-=1-2u-2(1-u)\Theta\le\beta
 \le1-2l+2l\Theta=b_+ .                                      \tag{1}
\]
There is no reflection of \(p\) to \(\min(p,1-p)\).

Expand the absolute cubic using the same conditional marginal:
\[
 \rho:=\E|X|^3
 =\ell^3\left\{pq(1-2pq)+3\E|Z-p|U^2
                    +\E[\operatorname{sgn}(Z-p)U^3]\right\}. \tag{2}
\]
The linear term vanishes conditionally.
Since \(|Z-p|\le\max(p,q)\), \(|U|^3\le(3/5)U^2\) and
\(\ell^3\le(pq)^{-3/2}\), the non-skeleton contribution is at most
\[
 \frac{3\max(p,q)+3/5}{(pq)^{3/2}}\,\frac{\lambda}{n}
 \le24\frac{\lambda}{n}.
\]
This coefficient is valid on the whole range, since \(pq\ge6/25\),
\(\max(p,q)\le3/5\), and exactly
\[
 24^2(6/25)^3-(12/5)^2=34416/15625>0.                       \tag{3}
\]
The function \((1-2v)/\sqrt v\) decreases for \(0<v\le1/4\);
the positive skeleton in (2) is multiplied by
\((1-\theta)^{3/2}\le1\).
If \(s_I\le\sqrt{v_-}\), this proves the usable same-law ceiling
\[
 \rho\le r_I:=\frac{1-2v_-}{s_I}+\frac{24(3/40)}{42}.         \tag{4}
\]
No separate incompatible moment maxima are asserted to constitute a law.

\paragraph{Contact and whole-ray exclusion.}
The retained fixed-target contact inequality is
\[
 0\le H_n(H,t)+c\{-|t|^3+\tfrac32\rho t^2+3\beta t+\rho e_n\}
       -na,\qquad t\in\operatorname{supp}P,
\]
with \(a\ge0\), \(c=C_{\rm E}\ge819/2000\) and
\(e_n<8/875\). The sharp predecessor is \(C_{n-1}\le c\).
Discarding \(-na\) preserves an upper bound.
For \(t>\sqrt3\), set
\[
 S=81/80,\quad A_t=(t^3-3t)/6,\quad H_4=t^4-6t^2+3,\quad
 E_t=t^5/4100+3t/820+17(3t^2-2)/78720.
\]
The retained uniform heat expansion and derivative estimate imply
\[
 H_n(H,\pm t)\le S^3\{\tfrac25(A_t+|H_4|/72)+E_t\},\qquad
 |\partial_t H_n(H,\pm t)|<\tfrac25(t^2/2+1/2+8/875).        \tag{5}
\]
These are bounds for all real raw \(H\). The absolute quartic term
avoids a positive-\(H_4\) assumption; \(t^2>3\) ensures the cubic sign
and the heat remainder's use of \(|t^2-1|=t^2-1\).

For \(r=r_I\), use \(b=-b_-\) on the negative ray and \(b=b_+\)
on the positive ray, and define
\[
 Q(t)=-t^3+\tfrac32rt^2+3bt+(8/875)r,\qquad
 L(t)=\tfrac25(t^2/2+1/2+8/875)+(819/2000)Q'(t).
\]
At a boundary check
\[
 P(t):=S^3\{\tfrac25(A_t+|H_4|/72)+E_t\}
                       +(819/2000)Q(t)<0,                 \tag{6}
\]
and \(Q,Q',Q'',L,L',L''<0\).
The sign of \(Q\) licenses replacing \(c\) by its lower bound in
(6). Since \(Q'''=-6\) and
\(L''=2/5-6(819/2000)<0\), these signs persist for all larger \(t\).
Equation (5), together with \(Q'<0\), then makes
\(H_n(H,\pm t)+cQ(t)\) strictly decreasing. Thus the entire exterior
ray, including its boundary, contains no support point.

Take the rational boundaries
\[
 t_-=E_I(l+1/2),\qquad t_+=E_I(3/2-u).                       \tag{7}
\]
The independent exact calculation uses directed square roots of precision
\(10^{-12}\), checks all thirteen closed intervals, and proves every
boundary satisfies \(t_\pm^2>3\).
All 26 inequalities (6) and all 156 continuation signs pass strictly.
The least favorable contact upper bound is on panel 40's positive ray
and is less than \(-.213\). Complete rational values are recorded in
the accompanying JSON output; no sampled positivity is used.

Consequently \(X\in(-t_-,t_+)\). Since \(\ell\ge E_I\) and
\(p\in[l,u]\), this implies \(Y=X/\ell+p\in(-1/2,3/2)\).
For \(Z=0\), the actual central gap gives
\(U=Y\le p-A/\ell=(\ell p-A)/\ell\le\delta<.072\).
For \(Z=1\), it gives
\(U=Y-1\ge-(\ell q-B)/\ell\ge-\delta>-.072\).
The exterior and inward bounds together yield \(|U|<1/2\).
At \(\lambda=0\), conditional centering and zero conditional variance
already give \(U=0\); no division by noise is used.
The shared probability and noise endpoints are included.

\paragraph{Scope and verification.}
The earlier moment proof supplies
\(Y\in(-.6,1.6)\) and \(\delta\le.0719539227\) on panels 40--52;
their inward consequences give precisely the prior \(|U|<.6\)
used above. Thus the new conclusion is not an input to its own proof.
Seven negative controls reject, including an actual standardized law
above \(p=1/2\) that refutes the shortcut \(b_+=1-2l+\Theta\).
The checker imports only the Python standard library.
A numerical PASS verifies the recorded rational implications, while
the contact, heat and acquired-strip premises retain their stated
analytic status. No larger-noise confinement or finite-base result is
claimed.
\paragraph{Continuous replacement of the thirteen new boundary rows.}
There is a shorter verification with common constants
\[
 \Theta_0=37/5000,\qquad E_0=249/125,\qquad r_0=221/200.
\]
Indeed \(5/677<\Theta_0\), \(E_0^2<4(1-\Theta_0)\), and
\((r_0-3/70)^2(6/25)>(13/25)^2\), with \(r_0>3/70\).
These prove the requisite theta, scale and third-moment bounds for
all \(p\in[.4,.6]\), provided the same prior strip and inward bounds
are supplied. The acquired entry application remains only
\(p\in[.4,.53]\).

For the negative ray set \(z=p\); for the positive ray set \(z=1-p\).
These are physical support rays for the original law, with no CDF
reflection. The pointwise signed beta bound is
\(b=2z-1+2(1-z)\Theta_0\). At the boundary
\(t=E_0(z+1/2)\), this is the affine function
\[
 b(t)=\frac{2(1-\Theta_0)}{E_0}t-2+3\Theta_0,\qquad
 t\in[t_0,t_1]=[1.7928,2.1912].
\]
The quartic \(H_4\) is increasing on this interval and negative at
\(t_1\), hence \(|H_4|=-H_4\). Substituting \(b(t)\) into (6)
therefore gives one explicit quintic \(\widehat P(t)\), with leading
coefficient \(S^3/4100>0\).

The third derivative \(\widehat P'''\) is a convex quadratic.
The independent rational check gives strictly negative values at both
endpoints, so \(\widehat P'''<0\) throughout. Consequently
\(\widehat P'\) is concave. Its two endpoint values are strictly
positive; concavity puts it above their positive chord. Thus
\(\widehat P\) increases, and
\(\widehat P(t_1)<-.1779<0\) proves every boundary inequality.
The proposed additional condition \(\widehat P''(t_0)<0\) is false:
its value is approximately \(+.0228525021\). It is not used in this
corrected concavity argument; the failed preliminary check is preserved.

For the continuation gates distinguish the boundary-family polynomial
\(\widehat Q(t)=Q_{r_0,b(t)}(t)\) from derivatives holding beta fixed:
\[
 R_1(t)=-3t^2+3r_0t+3b(t),\qquad R_2(t)=-6t+3r_0,
\]
\[
 \widehat L(t)=\tfrac25(t^2/2+1/2+8/875)+(819/2000)R_1(t),
 \quad L_1(t)=\tfrac25t+(819/2000)R_2(t),
 \quad L_2=\tfrac25-6(819/2000).
\]
The exact gates \(\widehat Q''(t_0)<0\) and
\(\widehat Q'(t_1)>0\) prove \(\widehat Q\) increases, while
\(\widehat Q(t_1)<0\). Both \(R_1\) and \(\widehat L\) have
negative second derivatives and negative first derivatives at \(t_0\),
so they decrease from strictly negative values there.
The affine functions \(R_2,L_1\) also decrease from negative values,
and \(L_2<0\). All six fixed-beta radial signs therefore hold at
every boundary. Their radial continuation was already justified above.
This proves the same strict half-amplitude result without a new
probability grid. The acquired thirteen-panel strip remains a separate
antecedent; it has not been removed.
\end{document}
